\documentclass[10.5pt,a4paper]{article}

\usepackage[T1]{fontenc}
\usepackage[utf8]{inputenc}
\usepackage{ebgaramond}
\usepackage{fontspec}
\newfontfamily\kr{Noto Serif CJK KR}
\usepackage[margin=0.9in,top=0.8in,bottom=0.7in]{geometry}
\usepackage{titlesec}
\usepackage{enumitem}
\usepackage{tabularx}
\usepackage{xcolor}
\usepackage[hidelinks]{hyperref}

\definecolor{accent}{HTML}{C0432A}
\definecolor{muted}{HTML}{6B6560}
\definecolor{ink}{HTML}{1A1916}

\hypersetup{colorlinks=true, urlcolor=accent, linkcolor=accent}
\color{ink}
\pagestyle{empty}
\setlength{\parindent}{0pt}

% ---- Section styling ----
\titleformat{\section}
  {\normalfont\large\scshape\color{accent}}{}{0em}{}[\vspace{-0.6em}{\color{muted!40}\titlerule}\vspace{0.2em}]
\titlespacing{\section}{0pt}{7pt}{3pt}

% ---- Date / content row (minipages so \\ line-breaks inside the content) ----
\newcommand{\entry}[2]{%
  \noindent
  \begin{minipage}[t]{0.16\linewidth}\raggedright{\small\color{muted}#1}\end{minipage}%
  \hspace{0.02\linewidth}%
  \begin{minipage}[t]{0.80\linewidth}\raggedright #2\end{minipage}%
  \par\vspace{3.5pt}%
}
% ---- Skills row ----
\newcommand{\skill}[2]{%
  \noindent
  \begin{minipage}[t]{0.22\linewidth}\raggedright{\color{muted}#1}\end{minipage}%
  \hspace{0.01\linewidth}%
  \begin{minipage}[t]{0.75\linewidth}\raggedright #2\end{minipage}%
  \par\vspace{3pt}%
}

\newcommand{\me}[1]{\textbf{\underline{#1}}}

\begin{document}

% ======================= HEADER =======================
{\LARGE\bfseries Myungjin Lee}\quad{\color{muted}\large \kr{이명진}}\\[2pt]
{\color{muted} M.E. in Computer Science \& Artificial Intelligence \textbullet\ Dongguk University}\\[1pt]
{\color{muted} AI for Drug Discovery \textbullet\ ADMET Prediction \textbullet\ Drug Repurposing \textbullet\ Graph Learning \textbullet\ XAI}\\[4pt]
{\small
\href{mailto:leemj637@gmail.com}{leemj637@gmail.com} \ \textbullet\ \
\href{https://github.com/myungjin00}{github.com/myungjin00} \ \textbullet\ \
\href{https://myungjin00.github.io}{myungjin00.github.io} \ \textbullet\ \
Seoul, South Korea}

% ======================= INTEREST =======================
\section{Research Interests}
\textbf{AI for Drug Discovery:} Molecular Graph Representation Learning, ADMET Property Prediction, Drug Repurposing, Drug--Target Interaction (DTI), Drug Response Prediction, Explainable AI (XAI).\\[2pt]
\textbf{ML \& Data Mining for Bioinformatics:} Multi-omics Data Integration, Network Biology, Differential Expression Analysis, Dimensionality Reduction, Multimodal AI for Biomedical Problems.

% ======================= EDUCATION =======================
\section{Education}
\entry{Mar 2024 -- Aug 2026}{\textbf{M.E. in Computer Science and Artificial Intelligence} \\ Dongguk University, Seoul, Korea \\ {\color{muted}Major: AI in Healthcare and Medicine \textbullet\ GPA 4.04/4.5}}
\entry{Mar 2023 -- Feb 2024}{\textbf{B.E. in Industrial \& Management Engineering} (Advanced Major Program) \\ Induk University, Seoul, Korea \\ {\color{muted}Convergence Minor in Deep Learning \textbullet\ GPA 4.43/4.5 \textbullet\ Valedictorian}}
\entry{Mar 2020 -- Feb 2023}{\textbf{A.E. in Industrial \& Management Engineering} \\ Induk University, Seoul, Korea \\ {\color{muted}GPA 4.41/4.5 (Major 4.42/4.5) \textbullet\ Valedictorian}}

% ======================= RESEARCH EXPERIENCE =======================
\section{Research Experience}
\entry{Mar 2024 -- Present}{\textbf{PRISM Lab}, Dongguk University, Seoul, Korea \\
{\color{muted}\textbf{Position:} Graduate Researcher (M.E.) --- Advisor: \href{https://scholar.google.com/citations?user=d19A738AAAAJ}{Dr.\ Sangsoo Lim}} \\
{\color{muted}\textbf{Role:} Molecular graph representation learning and explainable AI for drug discovery --- interpretable ADMET prediction from molecular substructures (first-author SJoINT; co-author MAGNET); a drug-repurposing model developed and registered to the national bio-data platform (KBDI); and bulk / single-cell RNA-seq omics analysis for target discovery.}}

% ======================= RESEARCH (MANUSCRIPTS) =======================
\section{Research}
{\small\color{muted}$^{*}$ Corresponding author; \underline{underlined} name is the CV owner.}\\[3pt]
\entry{2026}{\me{Lee M.}, Mo J., Kang M., Lim S.$^{*}$ \\ \textbf{SJoINT: Substructure-Driven Junction Tree for Interpretable ADMET Prediction} \\ {\color{muted}\textit{Manuscript in preparation.} Dual-view learning of atom-level graphs and substructure-level junction trees via structure-constrained bidirectional cross-attention with augmentation-free contrastive pretraining (ZINC250K); evaluated on MoleculeNet ADMET and MoleculeACE activity-cliff tasks.}}
\entry{2026}{Mo J., \me{Lee M.}, Lee S., Lim S.$^{*}$ \\ \textbf{MAGNET: Cross-view Molecular Graph Learning for Interpretable ADMET Prediction} \\ {\color{muted}\textit{Bioinformatics} (under review). Multi-view molecular graph framework integrating BRICS, junction-tree, and Murcko fragmentations via cross-view meta-graphs, with self-supervised pretraining on ZINC250K and ChemBERTa-enhanced fine-tuning across 10 MoleculeNet ADMET benchmarks.}}

% ======================= PRESENTATIONS =======================
\section{Conference Presentations}
\entry{Jun 2026}{\textbf{Structure-Constrained Bidirectional Cross-Attention for Self-Supervised Molecular Representation Learning} \\ {\color{muted}Poster, Korea Computer Congress (KCC) 2026, KIISE, Korea.}}
\entry{Oct 2025}{\textbf{Weighted Junction-Tree Nodes for Enhanced Interpretability in ADMET Tasks} \\ {\color{muted}Poster, BIOINFO 2025, KSBI, Korea.}}

% ======================= GRANTS =======================
\section{Grants}
\entry{2024 -- 2025}{\textbf{Master's Student Research Grant} (\kr{석사과정생 연구장려금지원사업}) --- \textbf{Principal Investigator} \\ {\color{muted}Ministry of Science and ICT (MSIT). Individual research grant on functional-group junction-tree-based multi-task molecular graph representation learning. Advisor: Dr.\ Sangsoo Lim.}}

% ======================= RESEARCH PROJECTS =======================
\section{Research Projects (Participating Researcher)}
\entry{2025 -- 2026}{\textbf{Multimodal AI-based Target Discovery and Drug Validation for Aging-related MASLD} \\ {\color{muted}Individual Basic Research Program, NRF (MSIT). PI: Dr.\ Sangsoo Lim; Collaborator: Dr.\ Gung Lee (Mayo Clinic, USA).} \\ {\color{muted}Preliminary omics/bioinformatics for target discovery: curated public \& experimental bulk/single-cell RNA-seq; QC $\rightarrow$ DEG $\rightarrow$ GSEA/ssGSEA $\rightarrow$ module scoring; characterized lipid-metabolism reprogramming and cellular-senescence signatures (SenMayo/SASP) in aging mouse liver, with a supporting cross-check in public human data; prioritized candidate genes and Plin2-defined senescent-cell subsets.}}
\entry{2025 -- 2026}{\textbf{Open AI Drug Discovery and Data Analysis Platform Center (Bio-SYNAPSE)} \\ {\color{muted}Bio \& Medical Technology Development Program, NRF (MSIT). PI: Dr.\ Minho Lee (host).} \\ {\color{muted}\textbf{Drug repurposing:} built a standardized PrimeKG-based benchmark to evaluate five models (TxGNN, DREAMWalk, DRHGCN, AMDGT, SVGA) across 190 diseases / 3{,}379 drugs; developed DC-VGAE and a performance-based ensemble. \textbf{ADMET:} benchmarked and developed an interpretable ADMET model using junction-tree / Murcko / BRICS substructures with cross-attention. \textbf{Deployment:} packaged the drug-repurposing model as a Docker image and registered it to the national bio-data platform (KBDI).}}

% ======================= COMPETITIONS =======================
\section{Competitions}
\entry{Aug 2024}{\textbf{2024 Samsung AI Challenge --- Machine Learning Force Fields} \quad {\color{accent}Top 10\% (11th)} \\ {\color{muted}Samsung Electronics SAIT $\cdot$ DACON. Built MLFF models predicting atomic energies and forces with uncertainty quantification for semiconductor-material simulation (team DAI).}}
\entry{2024}{\textbf{DREAM Olfactory Mixtures Prediction Challenge} \\ {\color{muted}International DREAM Challenge, Sage Bionetworks. Graph neural network predicting olfactory similarity of molecular mixtures.}}

% ======================= HONORS =======================
\section{Honors \& Awards}
\entry{2023}{\textbf{Highest Academic Achievement Award} \\ {\color{muted}Induk University, Korea.}}
\entry{2020 -- 2023}{\textbf{Merit-based Academic Scholarship} --- awarded every semester \\ {\color{muted}Induk University, Korea (incl.\ GRAPE Talent Scholarship).}}

% ======================= PATENTS =======================
\section{Patents \& Software Registrations}
\entry{Oct 2025}{An Embedding Method via Multi-view Graph Integration of Molecular Structures \\ {\color{muted}Software Registration, Korea Copyright Commission, No.\ C-2025-040207.}}
\entry{Jan 2025}{System and Method for Predicting Olfactory Characteristics of Odor Mixture Data \\ {\color{muted}Korean Patent Application (pending), No.\ 10-2025-0003146.}}
\entry{Nov 2024}{Deep Learning Model for Similarity Prediction of Graph-based Mixtures \\ {\color{muted}Software Registration, Korea Copyright Commission, No.\ C-2024-042597.}}

% ======================= SKILLS =======================
\section{Technical Skills}
\skill{Programming}{Python, R (basic), Java, C\#, JavaScript, SQL, HTML}
\skill{Deep Learning}{PyTorch, PyTorch Geometric, Hugging Face Transformers, scikit-learn}
\skill{Data Analysis \& Visualization}{pandas, NumPy, anndata, seaborn, matplotlib}
\skill{Cheminformatics}{RDKit, DeepChem, TorchDrug, ChemPy}
\skill{Bioinformatics}{Scanpy, Seurat, bulk \& single-cell RNA-seq, differential expression (DESeq2 / PyDESeq2, MAST), GSEA / ssGSEA / GSVA, GO enrichment}
\skill{Tools \& Databases}{Git, Linux, Docker, MySQL}

% ======================= LANGUAGES =======================
\section{Languages}
\skill{Korean}{Native}
\skill{English}{TOEIC 860 (LC 440 / RC 420)}

\end{document}
